\section{Source lock and literature boundary}
\label{sec:source}

\begin{definition}[Frozen object]
For all $n\ge2$, let $X_n=P^1(\ZZ/n\ZZ)$ and let $S,R$ act by
$S[a:b]=[-b:a]$ and $R[a:b]=[-b:a+b]$.  The global object is the union of
these blocks with inherited cusp correspondences between $n$ and $2n$, and
between $n$ and $3n$, at the distinguished cusp $[1:0]$.  The roof is the
Paper-32 roof and the marker $z$ still counts one inherited graph edge.
\end{definition}

The only new operation is a rational chain quotient and an explicit square
cell attachment.  The candidate is forbidden to consult $|X_n|-(n+1)$,
factorization, a prime table, an accepted-support table, target zeros, fitted
coefficients, or Route B.  Arithmetic strata in the experiment are
post-census labels only.

The closest mathematical collision is classical.  The quotient by
$1+S$ and $1+R+R^2$ on the projective line over $\ZZ/N\ZZ$ is the Manin-symbol
presentation of relative modular-symbol homology \citep{manin1972,merel1991}.
Ihara, Hashimoto, Bass, and Stark--Terras provide the surrounding graph-zeta
and nonbacktracking determinant vocabulary
\citep{ihara1966,hashimoto1989,bass1992,starkterras1996}.  Kac--Ward type
signed walk determinants require extra embedding or spin data
\citep{kacward1952,cimasoni2010}.  More recent graph-zeta work confirms that
weighted and alternating determinant technology is active, but it does not
change the present ownership test \citep{lau2025,ishikawa2026,lu2026}.
Recent explicit character formulae for congruence subgroups likewise reinforce
that character variation is level-wide structure, not a field-only selector
\citep{zhu2025}.

Thus the novelty claim is deliberately narrow.  We do not present Manin
relations, modular symbols, or graph zetas as new.  The contribution is the
closed-form no-go certificate for this specific Route-A source program:
killing the universal cells is possible, but it is generic topology, not
prime-selective recurrence.

\begin{figure}[t]
  \centering
  \begin{tikzpicture}[>=Latex,scale=1.0,
  sorbit/.style={circle,draw=formalblue,thick,fill=blue!8,minimum size=10mm,font=\small},
  rorbit/.style={circle,draw=deepgreen,thick,fill=green!8,minimum size=10mm,font=\small},
  edge/.style={line width=0.7pt,charcoal}]
  \node[sorbit] (S0) at (0,0) {$O_S$};
  \node[rorbit] (R0) at (4.0,0) {$O_R$};
  \node[sorbit] (S1) at (0,-2.0) {$O'_S$};
  \node[rorbit] (R1) at (4.0,-2.0) {$O'_R$};
  \draw[edge,formalblue,bend left=28] (S0) to (R0);
  \draw[edge,deepgreen,bend right=28] (S0) to (R0);
  \draw[edge] (S1) -- (R0);
  \draw[edge] (S0) -- (R1);
  \draw[edge,dashed] (S1) -- (R1);
  \node[align=center,font=\scriptsize] at (2.0,0.95)
    {$c=[1:0]$ and $Rc=[0:1]$ share both orbit endpoints};
  \node[align=center,font=\scriptsize] at (2.0,-2.75)
    {\textcolor{formalblue}{$e_c$} and \textcolor{deepgreen}{$e_{Rc}$}
    are parallel state edges; $e_c-e_{Rc}$ is the survivor};
\end{tikzpicture}

  \caption{The relation quotient is the cycle space of the incidence dessin
  whose vertices are $S$-orbits and $R$-orbits and whose edges are projective
  states.  The cusp state and its $R$ image form two parallel edges for every
  modulus.}
  \label{fig:dessin}
\end{figure}
